\chapter{Interpretation and Semantics}
\label{ch:semantics}

This chapter follows one small program from C++ to Lean, and each
section below re-examines it at a deeper level: first as text the gate
must recognize, then as values the evaluator computes, then as memory
the transfer theorem erases. The program is an accumulator exercised
twice, chosen because it is the smallest C++ shape that uses
construction, mutation, observation, and destruction together:

\begin{verbatim}
struct Acc {
  int32_t s;
  Acc() : s(0) {}
  ~Acc() {}
  void add(int32_t v) { s += v; }
  int32_t get() const { return s; }
};
int32_t acc_two(int32_t a, int32_t b) {
  Acc acc;
  acc.add(a);
  acc.add(b);
  return acc.get();
}
\end{verbatim}

Nothing about this program is heap-allocated, aliased, or clever, and
that is the point: it sits squarely inside the admitted subset, so
every stage of the pipeline has something to say about it.

\section{Scope and staging}

Interpretation proceeds in stages, and the stage boundary is also the
trust boundary. CIR text is parsed (trusted) into \texttt{RawIR}; oracle
verdicts about aliasing are read in (trusted); the verified gate
\texttt{validate} admits exactly the documented subset and lowers it to
\texttt{CoreIR}; the verified emitter turns \texttt{CoreIR} into Lean
definitions. Unverified text therefore never reaches the emitter without
passing the gate. Pointers do not survive as values across this boundary:
they become ownership roles (\texttt{owned}, \texttt{mutBorrow},
\texttt{sharedBorrow}), and the admitted shapes are enumerated normatively
in \texttt{docs/SUBSET.md}.

For our example, the compiler produces five functions, and the gate
reads them the way a bouncer reads IDs --- markers first, details
second. Trap and cleanup regions are rejected outright everywhere else
in the gate; this entry shape carries the explicit exemption, which is
why the scope below survives into \texttt{CoreIR}. Eliding allocas,
loads, stores, attributes, and location noise, the entry says:

\begin{verbatim}
cir.func ... @_Z7acc_twoii(%arg0: !s32i, %arg1: !s32i) -> !s32i {
  ...
  cir.call @_ZN3AccC2Ev(%3) : (...) -> ()
  cir.cleanup.scope {
    ...
    cir.call @_ZN3Acc3addEi(%3, %4) : (...) -> ()   // acc.add(a)
    ...
    cir.call @_ZN3Acc3addEi(%3, %5) : (...) -> ()   // acc.add(b)
    ...
    %6 = cir.call @_ZNK3Acc3getEv(%3) : (...) -> !s32i
    ...
    cir.return %7 : !s32i
  } cleanup normal {
    cir.call @_ZN3AccD2Ev(%3) : (...) -> ()         // ~Acc()
    ...
  }
  cir.trap
}
\end{verbatim}

What the gate keys on is not the arithmetic but the shape around it:
the constructor and destructor carry \texttt{cxx\_ctor} /
\texttt{cxx\_dtor} markers, the \texttt{this} pointers carry the
nonnull/dereferenceable/noundef attribute triple (never
\texttt{restrict}, which \texttt{this} cannot carry), the body is a
cleanup scope sequencing exactly five call sites, and the arithmetic
hides inside the callees. The single addition that matters textually is
one line in the \texttt{add} leaf:

\begin{verbatim}
%6 = cir.add nsw %5, %3 : !s32i
\end{verbatim}

The \texttt{nsw} flag is the difference between a checked operation and
a wrapping one, and the gate demands it for signed addition. Past the
gate, all of this becomes one canonical \texttt{CoreIR} function whose
body is a straight-line chain of calls, with the cleanup scope made
explicit sequencing:

\begin{verbatim}
.cleanup
  (.seq (.callRet "s0" "_ZN3AccC2Ev" [])
  (.seq (.callRet "s1" "_ZN3Acc3addEi" ["s0", "a"])
  (.seq (.callRet "s2" "_ZN3Acc3addEi" ["s1", "b"])
  (.seq (.callRet "s3" "_ZNK3Acc3getEv" ["s2"])
  (.seq (.callRet "u" "_ZN3AccD2Ev" ["s3"])
        (.return_ (.var "s3")))))))
\end{verbatim}

Methods have become functions over an explicit state word
(\texttt{s0} through \texttt{s3}); the object never appears. That
erasure, performed once per shape at admission time, is what makes the
next section possible. Two smaller erasures of the same kind predate
it and persist underneath: call composition is a program-level
\texttt{callRet} step over a DAG of call-free leaves (callers
evaluate by delegating to whole-function semantics, with errors
propagating past every later call); growth routines compose a level
higher, with composers calling proved composers through the depth-n
program step over frozen leaves, never by inlining; recursion bottoms
out in fuel too, with the single admitted recursive shape re-entering
itself on strictly smaller ranges through the same step and exhaustion
loud; and struct-by-value is a pair of
pure expressions --- field projection fusing member access with its
load, and point construction fusing field stores with the return ---
so neither callers nor small structs need memory to be stated.

\section{Value semantics}

The meaning of an admitted program is given by a loan-based, value-only
evaluation: environments map variables to values with borrow bookkeeping,
and there is no heap and no addresses. Failure inside the subset is made
explicit through a \texttt{Result} monad whose errors name the cause
(overflow, division by zero, out-of-bounds access, failed assertion,
uninitialized use). Loops are fuel-bounded, with fuel exhaustion reported
as an assertion failure --- a deliberate incompleteness that can never
become unsoundness, since the gate admits only loops paired with a length
bound the fuel covers. Loop bodies may exit early --- loop-scoped
\texttt{break} and \texttt{continue} with escape reported loudly,
plus early \texttt{return} short-circuiting the sequence ---
and switches arrive lowered to if-chains over exact constant pins.
Arithmetic is width-polymorphic over 32- and 64-bit words, with mixed
widths rejected loudly; sub-32-bit arithmetic is rejected too, because
the frontend already promotes it through 32-bit casts, so there is no
native small-width operation left to model. The value story has grown
past readers: a sequential writer to one address is admitted by
functionalization --- value in, updated value out --- and a disjoint
writer pair returns both updated words as a single pair value, the
writes commuting without sequencing. This evaluator is the
specification against which every emitted definition is proved correct,
one canonical function at a time.

Our accumulator's meaning is a three-line definition: start from zero,
thread two checked additions, read out the result. The constructor,
getter, and destructor are identities at value level (field-init to
zero, read-back, no-op), so only the additions can fail:

\begin{verbatim}
def accTwo (a b : BitVec 32) : Result (BitVec 32) :=
  match checkedAddI32 0 a with
  | .error e => .error e
  | .ok s1 => checkedAddI32 s1 b
\end{verbatim}

Evaluating \texttt{accTwo 3 4} walks the chain in the open. The
constructor yields \texttt{0}; the first call binds \texttt{s1} to the
result of \texttt{checkedAddI32 0 3}, which is \texttt{.ok 3}; the
second computes \texttt{checkedAddI32 3 4}, which is \texttt{.ok 7};
the getter returns \texttt{7} and the destructor passes it through, so
the whole program answers \texttt{.ok 7}. The \texttt{callRet} chain in
the \texttt{CoreIR} body mirrors this exactly: each \texttt{s}-variable
holds one intermediate word, and errors propagate past every later
call by construction. An overflowing input takes the other path ---
\texttt{checkedAddI32} reports \texttt{Overflow}, the match arm
delivers \texttt{.error}, and the emitted program agrees, which is
what the two error-path correctness lemmas pin down. Nothing was
allocated, freed, or framed in this telling; there was nowhere for a
memory to appear, because values were all the program ever touched.

\section{Memory model and transfer}

A second, addressful model exists for one purpose: to shrink trust in
the aliasing oracle. The memory model is a flat block map with
lightweight borrow tags, strong enough for exactly the project's
uniqueness sources (\texttt{restrict}/\texttt{noalias} attributes,
disjoint allocation results, length-paired access), and nothing more.
Per-shape footprint proofs plus a transfer theorem --- memory evaluation
coincides with value evaluation on admitted functions --- demote the
oracle verdicts from a soundness argument to a checked cache: CI asserts
agreement, but the proofs do not depend on belief. The recursive sort
is the exception that proves the pattern: its whole-function value
bridge closes the spec, so it ships no transfer file. The free
discipline follows the same incompleteness-over-unsoundness rule: a missing free
is tolerated as a leak, while a double free fails loudly through the
affine token and the gate. Call composition has a memory twin as well:
the program layer evaluates callers under the same fuel in both models,
so the transfer covers entries, not just leaves. Shapes whose
footprints depend on the caller (conditionally-live writers) were the
known residue of the purely equational story, prototyped natively over
the memory model; they are now admitted through an explicit framing
contract --- scalar stores land on pinned cells under the open contract,
and the transfer case-splits on the addresses.

Run the accumulator a third time, now with addresses switched on, and
watch nothing change. Each \texttt{this} pointer resolves through the
layout to a live block whose tag matches the allocating epoch; the
single-reference parameters bind plain values justified by the
attribute triple; the stores and loads inside the leaves hit their
cells and cross-check against the value-level reads. The transfer
theorem for this shape states the outcome: memory evaluation of the
entry agrees with \texttt{accTwo} on every input, ok and error paths
alike. The oracle's verdict for these two integer parameters ---
\texttt{unknown}, since integers never alias --- sits in the checked
cache, consulted and irrelevant at once. That is the whole moral of
the example, stated three ways across this chapter: the gate saw
markers and calls, the evaluator saw threaded words, and memory saw
tagged blocks, and all three told the same story about six lines of
C++.
